\section{What this project is about}
\label{sec:project}

Modern applications are built from dozens of small services calling each other. When
something goes wrong --- a database slows down, a service runs out of memory, a
co-tenant steals CPU --- the symptoms spread: many services get slow at once, and the
on-call engineer must work out \emph{which component is the culprit} and \emph{what
kind of fault it is}. That activity is called root-cause analysis (RCA).

Engineers do RCA by reading \emph{observability data}: numeric metrics (CPU, memory,
network per container), application logs, and distributed traces (records of each
user request as it hops between services). Our project adds a fourth, unusual source:
\emph{kernel traces} --- a complete record from the operating system of every system
call, scheduling decision, disk and network event, for every service, with no code
changes required. Kernel data can see things the other sources cannot (for example,
\emph{why} a process was waiting), but it is enormous and hard to read, which is
exactly where an LLM might help.

The project asks four questions:
\begin{enumerate}
  \item \textbf{Capability.} Can an LLM agent, given tools to query this data, match
  or beat existing automated RCA methods?
  \item \textbf{Robustness.} How much can the monitoring data be thinned --- fewer
  traces, coarser metrics, filtered logs --- before diagnosis quality breaks?
  \item \textbf{The value of kernel data.} What does the deep kernel layer actually
  contribute, and in what form is it useful?
  \item \textbf{Minimum budget.} What is the cheapest monitoring setup that still
  supports accurate diagnosis? (Observability is expensive to collect and store; if
  5\% of it is enough, that matters in practice.)
\end{enumerate}

A companion engineering track builds the agent itself as a product-shaped system ---
with a context builder that prepares an evidence briefing, and a library of reusable
investigation guides (``skills'') that anyone can extend --- aligned with the
industrial architecture our collaborators use.
